\ifdefined\gkver\else\def\gkver{student}\fi
\documentclass[\gkver]{gaokaozhenti}
\newcommand{\ccnum}[1]{{\fontspec{Noto Sans CJK JP}#1}}
\begin{document}

\chapter{2015年四川理综}
%% paperType: 理综物理部分

\begin{enumerate}
\item
%% number: 1
%% paperName: 2015年四川理综第 1 题 6 分
%% typeId: 1
%% type: 单选题
%% chapter: 功和能
%% point: 机械能守恒定律
%% method: 无
%% score: 6
%% degree: 650
%% duplicateId: 0
%% body:
在同一位置以相同的速率把三个小球分别沿水平、斜向上、斜向下方向抛出，不计空气阻力，则落在同一水平地面时的速度大小\xzanswer{A}

\fourchoices[answer=A]
{一样大}
{水平抛的最大}
{斜向上抛的最大}
{斜向下抛的最大}

%% answer: A
%% memo:
\memoanswer{
	根据动能定理：$mgh=\frac{1}{2}mv^{2}-\frac{1}{2}mv_{0}^{2}$，得落地速度 $v=\sqrt{v_{0}^{2}+2gh}$，而初速率 $v_{0}$ 相同，高度 $h$ 相同，故落地速度相同。选项 A 正确。
}

\item
%% number: 2
%% paperName: 2015年四川理综第 2 题 6 分
%% typeId: 1
%% type: 单选题
%% chapter: 机械波
%% point: 波长频率波速
%% method: 无
%% score: 6
%% degree: 650
%% duplicateId: 0
%% body:
平静湖面传播着一列水面波（横波），在波的传播方向上有相距 $3\Um$ 的甲、乙两小木块随波上下运动，测得两小木块每分钟上下 30 次，甲在波谷时，乙在波峰，且两木块之间有一个波峰。这列水面波\xzanswer{C}

\fourchoices[answer=C]
{频率是 $30\UHz$}
{波长是 $3\Um$}
{波速是 $1\Ums$}
{周期是 $0.1\Us$}

%% answer: C
%% memo:
\memoanswer{
	由两小木块每分钟上下 30 次知甲、乙的周期相同，为 $T=\frac{60}{30}\Us=2\Us$，频率 $f=0.5\UHz$，A、D 错误；甲在波谷时，乙在波峰，且两木块之间有一个波峰，则两木块间的距离为 $\frac{3}{2}\lambda=3\Um$，即波长 $\lambda=2\Um$，B 错误；波速 $v=\frac{\lambda}{T}=1\Ums$，C 正确，选项 C 正确。
}

\item
%% number: 3
%% paperName: 2015年四川理综第 3 题 6 分
%% typeId: 1
%% type: 单选题
%% chapter: 光学
%% point: 光的折射
%% method: 无
%% score: 6
%% degree: 650
%% duplicateId: 0
%% body:
直线 $P_{1}P_{2}$ 过均匀玻璃球球心 O，细光束 a、b 平行且关于 $P_{1}P_{2}$ 对称，由空气射入玻璃球的光路如图。a、b 光相比\xzanswer{C}

\onepicture[h=2.8cm]{figs/03.png}

\fourchoices[answer=C]
{玻璃对 a 光的折射率较大}
{玻璃对 a 光的临界角较小}
{b 光在玻璃中的传播速度较小}
{b 光在玻璃中的传播时间较短}

%% answer: C
%% memo:
\memoanswer{
	由于细光束 a、b 平行且关于 $P_{1}P_{2}$ 对称，射入玻璃球的入射角相等，由题图几何关系可知光线 a 的折射角小于 b 的折射角，根据折射定律 $\frac{\sin i}{\sin\gamma}=n$，玻璃对 b 光的折射率大，A 错误；由发生全反射的临界角 $\sin C=\frac{1}{n}$ 可得：玻璃对 b 光的临界角较小，B 错误；根据 $v=\frac{c}{n}$ 知 b 光在玻璃中的传播速度较小，C 正确；由于在玻璃中 b 的光程 d 比 a 大，由 $t=d/v$ 知 b 光在玻璃中的传播时间较长，D 错误，选项 C 正确。
}

\item
%% number: 4
%% paperName: 2015年四川理综第 4 题 6 分
%% typeId: 1
%% type: 单选题
%% chapter: 交流电
%% point: 交流电
%% method: 无
%% score: 6
%% degree: 650
%% duplicateId: 0
%% body:
小型发电机线圈共 $N$ 匝，每匝可简化为矩形线圈 abcd，磁极间的磁场视为匀强磁场，方向垂直于线圈中心轴 OO$'$，线圈绕 OO$'$ 匀速转动，如图所示。矩形线圈 ab 边和 cd 边产生的感应电动势的最大值都为 $e_{0}$，不计线圈电阻，则发电机输出电压\xzanswer{D}

\onepicture[h=3.4cm]{figs/04.png}

\fourchoices[answer=D]
{峰值是 $e_{0}$}
{峰值是 $2e_{0}$}
{有效值是 $\frac{\sqrt{2}}{2}Ne_{0}$}
{有效值是 $\sqrt{2}Ne_{0}$}

%% answer: D
%% memo:
\memoanswer{
	线圈 ab 边和 cd 边产生的感应电动势的最大值都为 $e_{0}$，单匝线圈总电动势最大值为 $2e_{0}$，由发电机的线圈为 $N$ 匝，故发电机的电动势的最大值 $E_{m}=2Ne_{0}$，其有效值为 $E=\frac{E_{m}}{\sqrt{2}}=\sqrt{2}Ne_{0}$，选项 D 正确。
}

\item
%% number: 5
%% paperName: 2015年四川理综第 5 题 6 分
%% typeId: 1
%% type: 单选题
%% chapter: 曲线运动
%% point: 人造地球卫星
%% method: 无
%% score: 6
%% degree: 450
%% duplicateId: 0
%% body:
登上火星是人类的梦想，“嫦娥之父”欧阳自远透露：中国计划于 2020 年登陆火星。地球和火星的公转视为匀速圆周运动，忽略行星自转影响。根据下表，火星和地球相比\xzanswer{B}

\begin{center}
\begin{tabular}{|c|c|c|c|}
\hline
行星 & 半径/$\Um$ & 质量/$\Ukg$ & 轨道半径/$\Um$ \\
\hline
地球 & $6.4\times10^{6}$ & $6.0\times10^{24}$ & $1.5\times10^{11}$ \\
\hline
火星 & $3.4\times10^{6}$ & $6.4\times10^{23}$ & $2.3\times10^{11}$ \\
\hline
\end{tabular}
\end{center}

\fourchoices[answer=B]
{火星的公转周期较小}
{火星做圆周运动的加速度较小}
{火星表面的重力加速度较大}
{火星的第一宇宙速度较大}

%% answer: B
%% memo:
\memoanswer{
	太阳对行星的万有引力提供行星绕太阳运动的向心力，设太阳的质量为 $M$，行星的质量为 $m$，行星运动的轨道半径为 $r$，则有 $G\frac{Mm}{r^{2}}=ma_{n}=m\frac{4\pi^{2}}{T^{2}}r$，行星公转周期 $T=2\pi\sqrt{\frac{r^{3}}{GM}}$，做圆周运动的加速度 $a_{n}=G\frac{M}{r^{2}}$，比较表中数据，A 错误，B 正确；设行星的半径为 $R$，行星上一个物体质量为 $m_{0}$，则 $G\frac{mm_{0}}{R^{2}}=m_{0}g$，行星表面的重力加速度 $g=G\frac{m}{R^{2}}$，故 $\frac{g_{\text{地}}}{g_{\text{火}}}=\frac{m_{\text{地}}R_{\text{火}}^{2}}{m_{\text{火}}R_{\text{地}}^{2}}=\frac{6.0\times10^{24}\times(3.4\times10^{6})^{2}}{6.4\times10^{23}\times(6.4\times10^{6})^{2}}=2.645$，火星表面的重力加速度较小，C 错误；设行星的第一宇宙速度为 $v$，则 $G\frac{mm_{0}}{R^{2}}=m_{0}\frac{v^{2}}{R}$，第一宇宙速度 $v=\sqrt{\frac{Gm}{R}}$，故 $\frac{v_{\text{地}}}{v_{\text{火}}}=\sqrt{\frac{m_{\text{地}}R_{\text{火}}}{m_{\text{火}}R_{\text{地}}}}=\sqrt{\frac{6.0\times10^{24}\times3.4\times10^{6}}{6.4\times10^{23}\times6.4\times10^{6}}}=2.23$，火星的较小，D 错误，选项 B 正确。
}

\item
%% number: 6
%% paperName: 2015年四川理综第 6 题 6 分
%% typeId: 2
%% type: 多选题
%% chapter: 电场
%% point: 带电粒子的圆周运动
%% method: 等效替代
%% score: 6
%% degree: 450
%% duplicateId: 0
%% body:
如图所示，半圆槽光滑、绝缘、固定，圆心是 O，最低点是 P，直径 MN 水平，a、b 是两个完全相同的带正电小球（视为点电荷），b 固定在 M 点，a 从 N 点静止释放，沿半圆槽运动经过 P 点到达某点 Q（图中未画出）时速度为零。则小球 a\xzanswer{BC}

\onepicture[h=3.4cm]{figs/06.png}

\fourchoices[answer=BC]
{从 N 到 Q 的过程中，重力与库仑力的合力先增大后减小}
{从 N 到 P 的过程中，速率先增大后减小}
{从 N 到 Q 的过程中，电势能一直增加}
{从 P 到 Q 的过程中，动能减少量小于电势能增加量}

%% answer: BC
%% memo:
\memoanswer{
	a 球从 N 点静止释放后，受重力 $mg$、b 球的库仑斥力 $F$ 和槽的弹力 $F_{N}$ 作用，a 球在从 N 到 Q 的过程中，$mg$ 与 $F$ 的夹角 $\theta$ 逐渐减小，若 $F$ 的大小不变，随着 $\theta$ 的减小 $mg$ 与 $F$ 的合力将逐渐增大；况且，由库仑定律和图中几何关系可知，随着 $\theta$ 的减小，$F$ 逐渐增大，因此合力一直增加，A 错误；在 a 球从 N 到 Q 的过程中，a、b 两小球距离逐渐变小，电场力做负功，电势能增加，C 正确；从 P 到 Q 的过程中，根据能的转化与守恒可知，其电势能增加量等于其机械能的减少量，a 球在 Q 点时的重力势能大于其在 P 点时的重力势能，因此该过程中动能一定在减少，且其减少量一定等于其电势能与重力势能增加量之和，D 错误；在从 P 到 Q 的过程中，a 球的动能在减少，其速率必减小，而开始在 N 点时速率为 0，开始向下运动段中，其速率必先增大，B 正确。选项 BC 正确。
}

\item
%% number: 7
%% paperName: 2015年四川理综第 7 题 6 分
%% typeId: 2
%% type: 多选题
%% chapter: 磁场
%% point: 带电粒子在磁场中的运动
%% method: 无
%% score: 6
%% degree: 250
%% duplicateId: 0
%% body:
如图所示，S 处有一电子源，可向纸面内任意方向发射电子，平板 MN 垂直于纸面，在纸面内的长度 $L=9.1\Ucm$，中点 O 与 S 间的距离 $d=4.55\Ucm$，MN 与 SO 直线的夹角为 $\theta$，板所在平面有电子源的一侧区域有方向垂直于纸面向外的匀强磁场，磁感应强度 $B=2.0\times10^{-4}\UT$。电子质量 $m=9.1\times10^{-31}\Ukg$，电量 $e=1.6\times10^{-19}\UC$，不计电子重力。电子源发射速度 $v=1.6\times10^{6}\Ums$ 的一个电子，该电子打在板上可能位置的区域的长度为 $l$，则\xzanswer{AD}

\onepicture[h=3.6cm]{figs/07.png}

\fourchoices[answer=AD]
{$\theta=90^\circ$ 时，$l=9.1\Ucm$}
{$\theta=60^\circ$ 时，$l=9.1\Ucm$}
{$\theta=45^\circ$ 时，$l=4.55\Ucm$}
{$\theta=30^\circ$ 时，$l=4.55\Ucm$}

%% answer: AD
%% memo:
\memoanswer{
	由于 $evB=m\frac{v^{2}}{r}$，则电子做匀速圆周运动的轨道半径 $r=\frac{mv}{eB}=4.55\times10^{-2}\Um=4.55\Ucm$，而 $d=4.55\Ucm$、$L=9.1\Ucm$，故 $r=d=\frac{L}{2}$。S 处的电子源可向纸面内任意方向发射电子，所以电子的轨迹是经过 S 的一系列半径为 $r$ 的等大圆，电子打在板上可能位置区域的长度为 $l=NA$（如下左图）。

	\onepicture[h=3.4cm]{figs/07_an1.png}

	由于 MN 与 SO 直线的夹角 $\theta$ 不定，要使电子的轨迹与 MN 相切，由几何关系可知，电子的轨迹圆心 C 一定在与 MN 距离为 $r$ 的平行线上，如右图所示。$l=4.55\Ucm$ 时，A 点与 O 点重合，电子运动轨迹线为 $\widehat{S_{1}A_{1}}$，由几何关系得 $S_{1}O$ 与 MN 的夹角 $\theta=30^\circ$，C 错误，D 正确；$l=9.1\Ucm$ 时，A 点与 M 点重合，电子运动轨迹线为 $\widehat{S_{2}A_{2}}$，由几何关系得 $S_{2}O$ 与 MN 的夹角 $\theta=90^\circ$，A 正确，B 错误，选项 AD 正确。

	\onepicture[h=2.8cm]{figs/07_an2.png}
}

\item
%% number: 8
%% paperName: 2015年四川理综第 8 题 11 分
%% typeId: 4
%% type: 实验题
%% chapter: 电路
%% point: 测电动势与内阻
%% method: 无
%% score: 11
%% degree: 450
%% duplicateId: 0
%% body:
（1）（6 分）某同学在“探究弹力和弹簧伸长的关系”时，安装好实验装置，让刻度尺零刻度与弹簧上端平齐，在弹簧下端挂 1 个钩码，静止时弹簧长度为 $l_{1}$。如图 1 所示，图 2 是此时固定在弹簧挂钩上的指针在刻度尺（最小分度是 1 毫米）上位置的放大图，示数 $l_{1}=$\tkanswer[1.8cm]{25.85}$\Ucm$。在弹簧下端分别挂 2 个、3 个、4 个、5 个相同钩码，静止时弹簧长度分别是 $l_{2}$、$l_{3}$、$l_{4}$、$l_{5}$。已知每个钩码质量是 $50\Ug$，挂 2 个钩码时，弹簧弹力 $F_{2}=$\tkanswer[1.8cm]{0.98}$\UN$（当地重力加速度 $g=9.8\Umsq$），要得到弹簧伸长量 $x$，还需要测量的是\tkanswer[2.2cm]{弹簧的原长}。作出 $F-x$ 曲线，得到弹力与弹簧伸长量的关系。

\twopicture[showlabel=false, hA=5.0cm, hB=3.6cm]
{figs/08a.png}
{figs/08b.png}

（2）（11 分）用实验测一电池的内阻 $r$ 和一待测电阻的阻值 $R_{x}$。已知电池的电动势约 $6\UV$，电池内阻和待测电阻阻值都为数十欧。可选用的实验器材有：

电流表 $A_{1}$（量程 $0\sim30\UmA$）；

电流表 $A_{2}$（量程 $0\sim100\UmA$）；

电压表 V（量程 $0\sim6\UV$）；

滑动变阻器 $R_{1}$（阻值 $0\sim5\UO$）；

滑动变阻器 $R_{2}$（阻值 $0\sim300\UO$）；

开关 S 一个，导线若干条。

某同学的实验过程如下：

Ⅰ．设计如图 1 所示的电路图，正确连接电路。

Ⅱ．将 $R$ 的阻值调到最大，闭合开关，逐次调小 $R$ 的阻值，测出多组 $U$ 和 $I$ 的值，并记录。以 $U$ 为纵轴，$I$ 为横轴，得到如图 2 所示的图线。

Ⅲ．断开开关，将 $R_{x}$ 改接在 B、C 之间，A 与 B 直接相连，其他部分保持不变。重复Ⅱ的步骤，得到另一条 $U-I$ 图线，图线与横轴 $I$ 的交点坐标为（$I_{0}$，0），与纵轴 $U$ 的交点坐标为（0，$U_{0}$）。

回答下列问题：

\begin{steps}
	\item 电流表应选用\tkanswer[1.6cm]{$A_{2}$}，滑动变阻器应选用\tkanswer[1.6cm]{$R_{2}$}；
	\item 由图 2 的图线，得电源内阻 $r=$\tkanswer[1.4cm]{25}$\UO$；
	\item 用 $I_{0}$、$U_{0}$ 和 $r$ 表示待测电阻的关系式 $R_{x}=$\tkanswer[2.6cm]{$\frac{U_{0}}{I_{0}}-r$}，代入数值可得 $R_{x}$；
	\item 若电表为理想电表，$R_{x}$ 接在 B、C 之间与接在 A、B 之间，滑动变阻器滑片都从最大阻值位置调到某同一位置，两种情况相比，电流表示数变化范围\tkanswer[1.6cm]{相同}，电压表示数变化范围\tkanswer[1.6cm]{不同}。（选填“相同”或“不同”）
\end{steps}

\twopicture[showlabel=false, hA=3.4cm, hB=4.2cm]
{figs/08c.png}
{figs/08d.png}

%% answer:
\jdanswer{
	（1）$25.85\Ucm$，$0.98\UN$，弹簧的原长；

	（2）①$A_{2}$，$R_{2}$；②$25\UO$；③$\frac{U_{0}}{I_{0}}-r$；④相同，不同。
}
%% memo:
\memoanswer{
	（1）由题图 2 指针指示，$l_{1}=25.85\Ucm$；挂 2 个钩码时，弹簧弹力 $F=2mg=2\times50\times10^{-3}\Ukg\times9.8\Umsq=0.98\UN$；弹簧伸长量 $x=l-l_{0}$，要得到弹簧伸长量 $x$，还需测量弹簧的原长 $l_{0}$。

	（2）①由于电池的电动势约 $6\UV$，电池内阻和待测电阻阻值都为数十欧，可知电路中的电流几十毫安，故电流表应选用 $A_{2}$；为了电路能正常工作、方便调节，多测几组数据，滑动变阻器应选用 $R_{2}$。②由图 1 电路可知，电压表测量的是路端电压，电流表测电路的总电流，故图 2 中，图线斜率绝对值即为电源的内阻，有：$r=|k|=\frac{5.5-4.0}{(80-20)\times10^{-3}}\UO=25\UO$。③当改接电路后，图线的斜率为电源的内阻与电阻 $R_{x}$ 之和，因此有：$r+R_{x}=|k'|=\frac{U_{0}-0}{I_{0}-0}$，故 $R_{x}=\frac{U_{0}}{I_{0}}-r$。④若电表为理想电表，$R_{x}$ 接在 B、C 之间与接在 A、B 之间，电路的总电阻可变范围不变，因此电流表的示数变化范围相同；$R_{x}$ 接在 B、C 之间时，电压表测滑动变阻器两端电压，而 $R_{x}$ 接在 A、B 之间时，电压表测滑动变阻器与 $R_{x}$ 两端电压的和，当滑动变阻器阻值为某一值时，电路的电流相同，而电压表的读数不同，故电压表示数变化范围不同。
}

\item
%% number: 9
%% paperName: 2015年四川理综第 9 题 15 分
%% typeId: 6
%% type: 计算题
%% chapter: 功和能
%% point: 动能定理
%% method: 无
%% score: 15
%% degree: 650
%% duplicateId: 0
%% body:
严重的雾霾天气，对国计民生已造成了严重的影响，汽车尾气是形成雾霾的重要污染源，“铁腕治污”已成为国家的工作重点，地铁列车可实现零排放，大力发展地铁，可以大大减少燃油公交车的使用，减少汽车尾气排放。

若一地铁列车从甲站由静止启动后做直线运动，先匀加速运动 $20\Us$ 达到最高速度 $72\Ukmh$，再匀速运动 $80\Us$，接着匀减速运动 $15\Us$ 到达乙站停住。设列车在匀加速运动阶段牵引力为 $1\times10^{6}\UN$，匀速阶段牵引力的功率为 $6\times10^{3}\UkW$，忽略匀减速运动阶段牵引力所做的功。

\onepicture[align=right, h=3.2cm]{figs/09.png}

\begin{enumerate}
	\item
	求甲站到乙站的距离；
	\item
	如果燃油公交车运行中做的功与该列车从甲站到乙站牵引力做的功相同，求公交车排放气体污染物的质量。（燃油公交车每做 1 焦耳功排放气体污染物 $3\times10^{-6}\Ug$）
\end{enumerate}

%% answer:
\jdanswer{
\begin{enumerate}
	\item
	$s=1950\Um$
	\item
	$M=2.04\Ukg$
\end{enumerate}
}
%% memo:
\memoanswer{
	（1）设列车匀加速直线运动阶段所用的时间为 $t_{1}$，距离为 $s_{1}$；在匀速直线运动阶段所用的时间为 $t_{2}$，距离为 $s_{2}$，速度为 $v$；在匀减速运动阶段所用的时间为 $t_{3}$，距离为 $s_{3}$；甲站到乙站的距离为 $s$。则

	$s_{1}=\frac{1}{2}a t_{1}^{2}$\ccnum{①}，$s_{2}=vt_{2}$\ccnum{②}，

	$s_{3}=\frac{v}{2}t_{3}$\ccnum{③}，$s=s_{1}+s_{2}+s_{3}$\ccnum{④}，

	由\ccnum{①}\ccnum{②}\ccnum{③}\ccnum{④}式联立，并代入数据解得 $s=1950\Um$\ccnum{⑤}。

	（2）设列车在匀加速直线运动阶段的牵引力为 $F$，所做的功为 $W_{1}$；在匀速直线运动阶段的牵引力的功率为 $P$，所做的功为 $W_{2}$。设燃油公交车做与该列车从甲站到乙站相同的功 $W$，将排放气态污染物质量为 $M$。则

	$W_{1}=Fs_{1}$\ccnum{⑥}，$W_{2}=Pt_{2}$\ccnum{⑦}，

	$W=W_{1}+W_{2}$\ccnum{⑧}，$M=(3\times10^{-9}\Ukg\cdot\UJ^{-1})\cdot W$\ccnum{⑨}，

	联立\ccnum{①}\ccnum{⑥}\ccnum{⑦}\ccnum{⑧}\ccnum{⑨}式并代入数据得 $M=2.04\Ukg$\ccnum{⑩}。
}

\item
%% number: 10
%% paperName: 2015年四川理综第 10 题 17 分
%% typeId: 6
%% type: 计算题
%% chapter: 电场
%% point: 带电粒子的直线运动
%% method: 无
%% score: 17
%% degree: 450
%% duplicateId: 0
%% body:
如图所示，粗糙、绝缘的直轨道 OB 固定在水平桌面上，B 端与桌面边缘对齐，A 是轨道上一点，过 A 点并垂直于轨道的竖直面右侧有大小 $E=1.5\times10^{6}\UNC$，方向水平向右的匀强电场。带负电的小物体 P 电荷量是 $2.0\times10^{-6}\UC$，质量 $m=0.25\Ukg$，与轨道间动摩擦因数 $\mu=0.4$，P 从 O 点由静止开始向右运动，经过 $0.55\Us$ 到达 A 点，到达 B 点时速度是 $5\Ums$，到达空间 D 点时速度与竖直方向的夹角为 $\alpha$，且 $\tan\alpha=1.2$。P 在整个运动过程中始终受到水平向右的某外力 $F$ 作用，$F$ 大小与 P 的速率 $v$ 的关系如表所示。P 视为质点，电荷量保持不变，忽略空气阻力，取 $g=10\Umsq$，求：

\onepicture[align=right, h=2.8cm]{figs/10.png}

\begin{center}
\begin{tabular}{|c|c|c|c|}
\hline
$v/\Ums$ & $0\leqslant v\leqslant 2$ & $2<v<5$ & $v\geqslant 5$ \\
\hline
$F/\UN$ & 2 & 6 & 3 \\
\hline
\end{tabular}
\end{center}

\begin{enumerate}
	\item
	小物体 P 从开始运动至速率为 $2\Ums$ 所用的时间；
	\item
	小物体 P 从 A 运动至 D 的过程，电场力做的功。
\end{enumerate}

%% answer:
\jdanswer{
\begin{enumerate}
	\item
	$\Delta t_{1}=0.5\Us$
	\item
	$W=-9.25\UJ$
\end{enumerate}
}
%% memo:
\memoanswer{
	（1）小物体 P 在速率从 0 至 $2\Ums$ 时，所受外力 $F_{1}=2\UN$，设其做匀变速直线运动的加速度为 $a_{1}$，经过时间 $\Delta t_{1}$ 速度为 $v_{1}$，则

	$F_{1}-\mu mg=ma_{1}$\ccnum{①}，$v_{1}=a_{1}\Delta t_{1}$\ccnum{②}，

	由\ccnum{①}\ccnum{②}式并代入数据得 $\Delta t_{1}=0.5\Us$\ccnum{③}。

	（2）小物体 P 从速率为 $2\Ums$ 运动至 A 点，受外力 $F_{2}=6\UN$，设其做匀变速直线运动的加速度为 $a_{2}$，则 $F_{2}-\mu mg=ma_{2}$\ccnum{④}。

	设小物体 P 从速度 $v_{1}$ 经过 $\Delta t_{2}$ 时间，在 A 点的速度为 $v_{2}$，则

	$\Delta t_{2}=0.55\Us-\Delta t_{1}$\ccnum{⑤}，$v_{2}=v_{1}+a_{2}\Delta t_{2}$\ccnum{⑥}。

	P 从 A 点至 B 点，受外力 $F_{2}=6\UN$、电场力和滑动摩擦力的作用，使其做匀变速直线运动，加速度为 $a_{3}$，电荷量为 $q$，在 B 点的速度为 $v_{3}$，从 A 点至 B 点的位移为 $x_{1}$，则

	$F_{2}-\mu mg-qE=ma_{3}$\ccnum{⑦}，$v_{B}^{2}-v_{2}^{2}=2a_{3}x_{1}$\ccnum{⑧}。

	P 以速度 $v_{3}$ 滑出轨道右端 B 点，设水平方向受外力为 $F_{3}$，电场力大小为 $F_{E}$，有

	$F_{E}=F_{3}$\ccnum{⑨}。

	$F_{3}$ 与 $F_{E}$ 大小相等方向相反，P 水平方向所受合力为零，所以 P 从 B 点开始做初速度为 $v_{3}$ 的平抛运动。设 P 从 B 点运动至 D 点用时为 $\Delta t_{3}$，水平位移为 $x_{2}$，由题意知

	$\frac{v_{3}}{g\Delta t_{3}}=\tan\alpha$\ccnum{⑩}，$x_{2}=v_{3}\Delta t_{3}$\ccnum{⑪}。

	设小物体 P 从 A 点至 D 点电场力做功为 $W$，则

	$W=-qE(x_{1}+x_{2})$\ccnum{⑫}。

	联立\ccnum{④}\ccnum{⑧}\ccnum{⑩}\ccnum{⑪}\ccnum{⑫}式并代入数据得

	$W=-9.25\UJ$\ccnum{⑬}。
}

\item
%% number: 11
%% paperName: 2015年四川理综第 11 题 19 分
%% typeId: 6
%% type: 计算题
%% chapter: 电磁感应
%% point: 电磁感应能量分析
%% method: 无
%% score: 19
%% degree: 250
%% duplicateId: 0
%% body:
如图所示，金属导轨 MNC 和 PQD，MN 与 PQ 平行且间距为 $L$，所在平面与水平面夹角为 $\alpha$，N、Q 连线与 MN 垂直，M、P 间接有阻值为 $R$ 的电阻；光滑直导轨 NC 和 QD 在同一水平面内，与 NQ 的夹角都为锐角 $\theta$。均匀金属棒 ab 和 ef 质量均为 $m$，长均为 $L$，ab 棒初始位置在水平导轨上与 NQ 重合；ef 棒垂直放在倾斜导轨上，与导轨间的动摩擦因数为 $\mu$（$\mu$ 较小），由导轨上的小立柱 1 和 2 阻挡而静止。空间有方向竖直的匀强磁场（图中未画出）。两金属棒与导轨保持良好接触。不计所有导轨和 ab 棒的电阻，ef 棒的阻值为 $R$，最大静摩擦力与滑动摩擦力大小相等，忽略感应电流产生的磁场，重力加速度为 $g$。

\onepicture[align=right, h=3.8cm]{figs/11.png}

\begin{enumerate}
	\item
	若磁感应强度大小为 $B$，给 ab 棒一个垂直于 NQ、水平向右的速度 $v_{1}$，在水平导轨上沿运动方向滑行一段距离后停止，ef 棒始终静止，求此过程 ef 棒上产生的热量；
	\item
	在（1）问过程中，ab 棒滑行距离为 $d$，求通过 ab 棒某横截面的电量；
	\item
	若 ab 棒以垂直于 NQ 的速度 $v_{2}$ 在水平导轨上向右匀速运动，并在 NQ 位置时取走小立柱 1 和 2，且运动过程中 ef 棒始终静止。求此状态下最强磁场的磁感应强度及此磁场下 ab 棒运动的最大距离。
\end{enumerate}

%% answer:
\jdanswer{
\begin{enumerate}
	\item
	$\frac{1}{4}mv_{1}^{2}$
	\item
	$\frac{2Bd(L-d\cot\theta)}{R}$
	\item
	$\frac{1}{L}\sqrt{\frac{mgR(\sin\alpha+\mu\cos\alpha)}{(\cos\alpha-\mu\sin\alpha)v_{2}}}$，$\frac{\mu L\tan\theta}{(1+\mu^{2})\sin\alpha\cos\alpha+\mu}$
\end{enumerate}
}
%% memo:
\memoanswer{
	（1）设 ab 棒的初动能为 $E_{k}$，ef 棒和电阻 $R$ 在此过程产生的热量分别为 $W$ 和 $W_{1}$，有

	$W+W_{1}=E_{k}$\ccnum{①}，$W=W_{1}$\ccnum{②}，

	由题有 $E_{k}=\frac{1}{2}mv_{1}^{2}$\ccnum{③}，得 $W=\frac{1}{4}mv_{1}^{2}$\ccnum{④}。

	（2）设在题设过程中，ab 棒滑行时间为 $\Delta t$，扫过的导轨间的面积为 $\Delta S$，通过 $\Delta S$ 的磁通量为 $\Delta\Phi$，ab 棒产生的电动势为 $E$，ab 棒中的电流为 $I$，通过 ab 棒某横截面的电量为 $q$，则

	$E=\frac{\Delta\Phi}{\Delta t}$\ccnum{⑤}，$\Delta\Phi=B\Delta S$\ccnum{⑥}，

	$I=\frac{q}{\Delta t}$\ccnum{⑦}，又 $I=\frac{2E}{R}$\ccnum{⑧}。

	\onepicture[h=3.2cm]{figs/11_an1.png}

	由图所示 $\Delta S=d(L-d\cot\theta)$\ccnum{⑨}。

	联立\ccnum{⑤}—\ccnum{⑨}，解得 $q=\frac{2Bd(L-d\cot\theta)}{R}$\ccnum{⑩}。

	（3）ab 棒滑行距离为 $x$ 时，ab 棒在导轨间的棒长 $L_{x}$ 为

	$L_{x}=L-2x\cot\theta$\ccnum{⑪}，此时 ab 棒产生的电动势 $E_{x}$ 为 $E_{x}=Bv_{2}L$\ccnum{⑫}，

	流过 ef 棒的电流 $I_{x}$ 为 $I_{x}=\frac{E_{x}}{R}$\ccnum{⑬}，ef 棒所受安培力 $F_{x}$ 为 $F_{x}=BI_{x}L$\ccnum{⑭}。

	联立\ccnum{⑪}—\ccnum{⑭}，解得 $F_{x}=\frac{B^{2}v_{2}L}{R}(L-2x\cot\theta)$\ccnum{⑮}。

	由\ccnum{⑮}式可得，$F_{x}$ 在 $x=0$ 和 $B$ 为最大值 $B_{m}$ 时有最大值 $F_{1}$。

	由题知，ab 棒所受安培力方向必水平向左，ef 棒所受安培力方向必水平向右，使 $F_{1}$ 为最大值的受力分析如图所示，图中 $f_{m}$ 为最大静摩擦力，有

	\onepicture[h=2.8cm]{figs/11_an2.png}

	$F_{1}\cos\alpha=mg\sin\alpha+\mu(mg\cos\alpha+F_{1}\sin\alpha)$\ccnum{⑯}。

	联立\ccnum{⑮}\ccnum{⑯}，得 $B_{m}=\frac{1}{L}\sqrt{\frac{mgR(\sin\alpha+\mu\cos\alpha)}{(\cos\alpha-\mu\sin\alpha)v_{2}}}$\ccnum{⑰}。

	\ccnum{⑰}式就是题目所求最强磁场的磁感应强度大小，该磁场方向可竖直向上，也可竖直向下。

	由\ccnum{⑮}式可知，$B$ 为 $B_{m}$ 时，$F_{x}$ 随 $x$ 增大而减小，$x$ 为最大 $x_{m}$ 时，$F_{x}$ 为最小值 $F_{2}$，由图可知

	$F_{2}\cos\alpha+\mu(mg\cos\alpha+F_{2}\sin\alpha)=mg\sin\alpha$\ccnum{⑱}。

	联立\ccnum{⑮}\ccnum{⑰}\ccnum{⑱}，得

	\onepicture[h=2.6cm]{figs/11_an3.png}

	$x_{m}=\frac{\mu L\tan\theta}{(1+\mu^{2})\sin\alpha\cos\alpha+\mu}$\ccnum{⑲}。
}

\end{enumerate}

\end{document}
